%document class
\documentclass[10pt,oneside]{book}
%%%% Page Info + Commands %%%%%{

%packages
\usepackage{pgfplots}
\pgfplotsset{compat=1.18}
\usepackage{amsmath}
\usepackage{amsfonts}
\usepackage{amsthm,letltxmacro}
\usepackage{amssymb}
\usepackage{enumerate}
\usepackage{enumitem}
\usepackage[dvipsnames]{xcolor}
\usepackage{changepage}

%This will shrink the page
\newcommand{\smallpage}{  \setlength \oddsidemargin{.25in}
            \setlength \textwidth{5.5in}
            \setlength \topmargin{0in}
            \setlength \textheight{9in}}


% Commands for sets of numbers
\newcommand{\Z}{\mathbb{Z}}\newcommand{\R}{\mathbb{R}}\newcommand{\N}{\mathbb{N}}

% gordie spacing and boxing commands
\newcommand{\gspc}{\vspace{0.13cm}} % 0.5 space
\newcommand{\gline}{\gspc\\} % 1.5 space
\newcommand{\gbline}{\gspc\gspc\\} % double space
\newcommand{\gbox}[1]{\fbox{\parbox{\linewidth}{#1}}} % multiline box command
\newcommand{\gmod}{\text{mod}\,}


\newcommand{\gimage}[3]{
\begin{figure}[h]   
    \centering          
    \includegraphics[width=#2\textwidth]{#1}  
    \caption{#3}
    \label{fig:#1}
\end{figure}\\
}

\newcommand{\centerit}[1]{{\begin{center} #1 \end{center}}}
\newcommand{\ul}[1]{\underline{#1}}

%%%%%%%% command for graphics %%%%%%%%%%%%%
\usepackage{fancyhdr}
%}

%%%% Page Setup %%%%%%%
\smallpage\sloppy
\pagestyle{fancy}
\lhead{\large{\textbf{Problem Set 8}}}
%\rfoot{\thepage/\pageref{LastPage} }
\setlength{\headheight}{10pt}


% Proof writing
\LetLtxMacro\oldproof\proof
\let\endoldproof\endproof
\renewenvironment{proof}[1][\proofname]
  {\vspace{0.2cm}\begin{adjustwidth}{2em}{2em}
   \oldproof[#1]}
  {\endoldproof
\end{adjustwidth}}

\begin{document}

\newcommand{\cg}[1]{\color{OliveGreen}#1\color{white}}
\newcommand{\cb}[1]{\color{BlueViolet}#1\color{white}}

\subsection*{Section 5.2}

\begin{enumerate}
	\item Use Fermat's theorem to verify that $17$ divides $11^{104}+1$.
	\begin{proof}
		We want to show that $17$ divides $11^{104}+1$. \gline
		Note that via the corollary to Fermat's Little Theorem that for any nonzero $a$ in a prime modulus, that $a^{17} = a$.\gline
		Consider that:
		\begin{align*}
			11^{104}+1 &= (11^{17})^6\cdot 11^2 + 1\\
			&\equiv 11^6\cdot 11^2 + 1\mod 17\\
			&\equiv 11^8 + 1 \mod 17
		\end{align*}
		Now, note that $11^2 = 121\equiv 2\gmod 17$. We can simplify our equation to be:
		\begin{align*}
			11^8 + 1 &\equiv (11^2)^4 + 1\mod 17\\
			&\equiv 2^4 + 1\mod 17\\
			&\equiv 16 + 1\mod 17\\
			&\equiv 0 \mod 17
		\end{align*}
		Hence, because $11^{104}+ 1 \equiv 0(\gmod 17)$, 17 divides $11^{104} + 1$, as desired. 
	\end{proof}
	\item \begin{enumerate}
		\item If $\gcd(a, 35)=1$, show that $a^{12}\equiv 1(\gmod 35)$.
		\begin{proof}
			We want to show that if $\gcd(a,35) = 1$, then $a^{12}\equiv 1 (\gmod 35)$. \gline
			Consider that $a$ and 35 are relatively prime, and thus share no common prime factors. Hence, we show that $a^{12}\equiv 1(\gmod 7)$ and $a^{12}\equiv 1(\gmod 5)$.
			\centerit{\ul{$a^{12}\equiv 1 (\gmod 7)$}}
			Via Fermat's little theorem, we have that $a^{6}\equiv 1 (\gmod 7)$. \\
			It follows that: 
			\begin{align*}
				a^{12}&=(a^{6})^2\\
				&\equiv (1)^2(\gmod 7)\\
				&\equiv 1 (\gmod 7)
			\end{align*}
			Hence, $a^{12}\equiv 1(\gmod 7)$.
			\centerit{\ul{$a^{12}\equiv 1 (\gmod 5)$}}
			Via Fermat's little theorem, we have that $a^{4}\equiv 1 (\gmod 5)$. \\
			It follows that: 
			\begin{align*}
				a^{12}&=(a^{4})^3\\
				&\equiv (1)^3(\gmod 5)\\
				&\equiv 1 (\gmod 5)
			\end{align*}
			Thus, $a^{12}\equiv 1(\gmod 5)$.
		\end{proof}
		Hence, because $a^{12}\equiv 1(\gmod 7)$ and $a^{12}\equiv 1(\gmod 5)$, it follows that \\$a^{12}\equiv 1(\gmod 35)$, as desired. 
		\item If $\gcd(a,42)\equiv 1$, show that $168 = 3\cdot 7 \cdot 8$ divides $a^6 - 1$. 
		\begin{proof}
			We want to show that if $\gcd(a,42)\equiv 1$, it follows that $168$ divides $a^6 - 1$.\gline
			We will first demonstrate that $a^6 \equiv 1(\gmod 42)$. \gline
			Consider that $a$ is relatively prime to $42 = 2\cdot 3\cdot 7$, so it shares no common prime factors. Thus, proving that $a^6\equiv 1(\gmod 2)$, $a^6\equiv 1(\gmod 3)$, and $a^6\equiv 1(\gmod 7)$ is equivalent to proving that $a^6\equiv 1(\gmod 42)$.
			\pagebreak
			\centerit{\ul{$a^6\equiv 1 (\gmod 2)$}}
			Because $a$ is odd, because it does not share the prime factor of 2 with 32, and any odd number to any power remains odd, thus $a\equiv 1 (\gmod 2)$ \\
			Hence $a^6\equiv 1 (\gmod 2)$. 
			\centerit{\ul{$a^6\equiv 1 (\gmod 3)$}}
			Consider that via Fermat's little theorem, $a^{3-1}=a^2\equiv 1(\gmod 3)$.\\
			It follows that $a^6 = (a^2)^3 \equiv (1)^3(\gmod 3) \equiv 1 (\gmod 3)$\\
			Hence $a^6\equiv 1 (\gmod 3)$. 
			\centerit{\ul{$a^6\equiv 1 (\gmod 7)$}}
			Consider that via Fermat's little theorem, $a^{7-1}=a^6\equiv 1(\gmod 7)$. No further proof needed.\gline
			Hence $a^6\equiv 1 (\gmod 42)$. \\\\
			Now, we show that $a^6 - 1\equiv 0 (\gmod 168)$ (equivalently $168\mid a^6 - 1$). \gline
			Because 168 prime factorizes into $2^3,3$, and 7, it remains that we must show that $a^6\equiv 1 (\gmod 8)$. \gline
			Consider that $a$ is an odd number because the prime factor of 2 does not divide it. Let $k\in \Z$ such that $a = 2k + 1$. It follows that:
			\begin{align*}
				a^6 &= (a^2)^3\\
				&= (4k^2 + 4k + 1)^3\\
				&= (4k(k+1) + 1)^3
			\end{align*}
			Becuase $k(k+1)$ is an even number (even times an odd), let $j$ such that $k(k+1) = 2j$. We then have that:
			\begin{align*}
				(4k(k+1)+1)^3 &= (8j + 1)^3
			\end{align*}
			Note that $(8j+1)^3$ is trivially factorable into the form $8l + 1, l\in\Z$, because an 8 is attached to each $j$ term. 
			\gline 
			Thus, we have that $a^6 \equiv 1 (\gmod 8)$. \gline
			Hence, because $a^6\equiv 1$ for mod 8, mod 3, mod 7, it follow that $a^6\equiv 1(\gmod 168)$, and thus, $a^6 - 1 \equiv 0 (\gmod 168)$, as desired. 
		\end{proof}

	\end{enumerate}
	\item Derive each fo the following congruences:
	\begin{enumerate}
		\item $a^{21}\equiv a(\gmod 15)$ for all $a$. 
		\begin{proof}
			Let $a$ be an integer. We want to show that $\forall a\in \Z$, we have that $a^{21}\equiv a(\gmod 15)$.\gline
			Via the corollary to Fermat's little theorem, we have that $a^{15}\equiv a (\gmod 15)$. We also have that $(a^3)^7\equiv a^7 (\gmod 3)\equiv (a^2)a \equiv a (\gmod 3)$ and $(a^5)^4a  \equiv a^5 \equiv a (\gmod 5)$.\gline
			Thus, we have that $a^{21} \equiv a (\gmod 15)$. 
		\end{proof}
	\end{enumerate}
	\item Exploy Fermat's theorem to prove that, if $p$ is an odd prime, then:
	\begin{enumerate}
		\item $1^{p-1}+2^{p-1}+3^{p-1}+\cdots+(p-1)^{p-1}\equiv-1(\gmod p)$
		\begin{proof}
			Let $p$ be an odd prime. We want to show that 
			\begin{align}
				1^{p-1}+2^{p-1}+3^{p-1}+\cdots+(p-1)^{p-1}\equiv-1(\gmod p).
			\end{align}
			Consider the representation of equation (1) in mod $p$ as:
			\begin{align*}
				1^{p-1}+2^{p-1}+\cdots + (-2)^{p-1} + (-1)^{p-1}(\gmod p)
			\end{align*}
			Because $p$ is an odd number, we write all of our values as positive:
			\begin{align*}
				1^{p-1}+2^{p-1}+\cdots + (2)^{p-1} + (1)^{p-1}(\gmod p)
			\end{align*}
			Via Fermat's little theorem, because $p$ is prime and all of these $a$ are positive, we rewrite our representation as:
			\begin{align*}
				\equiv 1+1+\cdots + 1 + 1(\gmod p)
			\end{align*}
			And because we have $p-1$ elements, we have that this equation is:
			\begin{align*}
				&\equiv p -1 (\gmod p)\\
				&\equiv -1 (\gmod p),
			\end{align*}
			as desired.
		\end{proof}
		\item $1^p + 2^p + 3^p + \cdots + (p-1)^p \equiv 0 (\gmod p)$
		\begin{proof}
			Let $p$ be an odd prime. We want to show that 
			\begin{align}
				1^p + 2^p + 3^p + \cdots + (p-1)^p \equiv 0 (\gmod p)
			\end{align}
			Consider via Fermat's theorem, that $a^p \equiv a(\gmod p)$. We rewrite equation (2) as:
			\begin{align*}
				1+2+3+\cdots + (p-1) \mod p
			\end{align*}
			Consider that for any $p-n$, $p-n\equiv -n(\gmod p)$. We then rewrite our equation as:
			\begin{align*}
				1 + 2+ 3 + \cdots + -3 + -2 + -1
			\end{align*}
			Because we have an even number of elements ($p-1$ elements , and $p$ is odd), we have that:
			\begin{align*}
				1 + 2+ 3 + \cdots + -3 + -2 + -1 \equiv 0 (\gmod p)
			\end{align*}
	Hence, equation (2) is true, as desired. 
		\end{proof}
	\end{enumerate}
\end{enumerate}


\end{document}
